% !TEX root = ../Franciszek_Zabierowski_mgr.tex
% Wspolne makra matematyczne (wartosci domyslne = notatki w University-Of-Physics-dev).
% Czcionka matematyczna zostaje CM (newtxsf w preambule jest zakomentowane).

% ---------- Re / Im: R i I jako \mathbb, "e" i "m" prostym krojem ----------
\DeclareRobustCommand{\Realis}{\mathop{\mathbb{R}\mathrm{e}}\nolimits}
\DeclareRobustCommand{\Imaginaris}{\mathop{\mathbb{I}\mathrm{m}}\nolimits}

% ---------- Powiekszone operatory ----------
% \LargeOp[skala]{operator}   sumy, iloczyny, \bigcup ...: granice nad/pod w trybie display, z boku w tekscie.
\newcommand{\LargeOp}[2][1.4]{\mathop{\vcenter{\hbox{\scalebox{#1}{$\displaystyle#2$}}}}\displaylimits}
\newcommand{\LargeSum}[1][1.5]{\LargeOp[#1]{\sum}}
\newcommand{\LargeProd}[1][1.5]{\LargeOp[#1]{\prod}}

% Calki: w trybie display granice NAD i POD znakiem, przesuniete jak w TeX-owej korekcji kursywy
% (gorna w prawo, dolna w lewo o polowe \fontcharic znaku), w tekscie zwykle granice z boku.
% Skladnia jak zawsze: \LargeInt[skala]_{a}^{b} f
\makeatletter
\newdimen\LI@dx
\newcommand{\LI@tilt}{0.25}% przesuniecie granic = \LI@tilt * skala * 1em (gorna w prawo, dolna w lewo o polowe)
\def\LI@go#1#2{\def\LI@s{#1}\def\LI@op{#2}\def\LI@lo{}\def\LI@hi{}\LI@parse}
\def\LI@parse{\@ifnextchar_{\LI@sub}{\@ifnextchar^{\LI@sup}{\LI@done}}}
\def\LI@sub_#1{\def\LI@lo{#1}\LI@parse}
\def\LI@sup^#1{\def\LI@hi{#1}\LI@parse}
\def\LI@text{\mathop{\LI@op}\nolimits_{\LI@lo}^{\LI@hi}}
\def\LI@display{%
  \LI@dx=\LI@s em \relax
  \LI@dx=\LI@tilt\LI@dx
  \mathop{\vcenter{\offinterlineskip\halign{\hfil##\hfil\cr
    \hbox{\kern\LI@dx$\scriptstyle\LI@hi$}\cr
    \noalign{\kern1pt}
    \hbox{\scalebox{\LI@s}{$\displaystyle\LI@op$}}\cr
    \noalign{\kern1pt}
    \hbox{\kern-\LI@dx$\scriptstyle\LI@lo$}\cr}}}\nolimits}
\def\LI@done{\mathchoice{\LI@display}{\LI@text}{\LI@text}{\LI@text}}

% \oiint: amsmath nie ma; dwie calki z kolkiem
\newcommand{\oiint}{\mathop{\vcenter{\hbox{\ooalign{\hfil$\displaystyle\iint$\hfil\cr\hfil$\displaystyle\bigcirc$\hfil\cr\vphantom{$\displaystyle\iint$}\cr}}}}\nolimits}

\newcommand{\LargeInt}[1][1.25]{\LI@go{#1}{\int}}
\newcommand{\LargeIInt}[1][1.25]{\LI@go{#1}{\iint}}
\newcommand{\LargeIIInt}[1][1.25]{\LI@go{#1}{\iiint}}
\newcommand{\LargeOint}[1][1.25]{\LI@go{#1}{\oint}}
\newcommand{\LargeOiint}[1][1.25]{\LI@go{#1}{\oiint}}
\makeatother
